% Ngân hàng bài tập — chương 2
% Chương được xác định bởi tên tệp này.

\begin{exo}[Từ nước đá lạnh đến nước sôi: ước lượng bậc độ lớn]{chauffage-glace-eau-ebullition}
\seoready{true}
\keywords{phép đo nhiệt lượng, chuyển pha, nóng chảy, hóa hơi, nhiệt ẩn, bậc độ lớn}

\textbf{Số liệu (ở áp suất khí quyển):}
\[
\begin{aligned}
c_{\text{ice}} &= 2{,}1\times10^3\,\text{J}\!\cdot\!\text{kg}^{-1}\!\cdot\!\text{K}^{-1},
& L_f &= 3{,}34\times10^5\,\text{J}\!\cdot\!\text{kg}^{-1},\\
c_{\text{water}} &= 4{,}18\times10^3\,\text{J}\!\cdot\!\text{kg}^{-1}\!\cdot\!\text{K}^{-1},
& L_v &= 2{,}26\times10^6\,\text{J}\!\cdot\!\text{kg}^{-1}.
\end{aligned}
\]
Ở áp suất khí quyển, một kilôgam nước đá ban đầu ở $-100\,{}^\circ\text{C}$ được đun nóng dần.
\begin{enumerate}
  \item Tính năng lượng cần để đưa nước đá từ $-100\,{}^\circ\text{C}$ lên $0\,{}^\circ\text{C}$.
  \item Tính năng lượng cần để làm tan hoàn toàn nước đá ở $0\,{}^\circ\text{C}$.
  \item Tính năng lượng cần để đun nước lỏng từ $0\,{}^\circ\text{C}$ lên $100\,{}^\circ\text{C}$.
  \item Tính năng lượng bổ sung cần để hóa hơi hoàn toàn nước ở $100\,{}^\circ\text{C}$.
  \item Giai đoạn nào cần nhiều năng lượng nhất? So sánh năng lượng để đun một kilôgam nước từ $0\,{}^\circ\text{C}$ lên $100\,{}^\circ\text{C}$ với thế năng của khối lượng $m$ rơi từ độ cao mười mét. Phải cho khối lượng bao nhiêu rơi để giải phóng cùng năng lượng?
  \item Xét quá trình ngược lại. Khối lượng $m_v$ hơi nước ngưng tụ trên kính chắn gió và nước lỏng tạo thành không nguội thêm. Biểu diễn nhiệt $Q_{\text{rel}}$ tỏa ra khi ngưng tụ và tính cho $m_v=1{,}0\times10^{-2}\,\text{kg}$. Ở nhiệt độ kính, lấy $L_v=2{,}45\times10^6\,\text{J}\!\cdot\!\text{kg}^{-1}$.
\end{enumerate}
\begin{indication}
Khi đun nóng mà không chuyển pha, dùng $Q=mc\Delta T$. Khi chuyển pha ở nhiệt độ không đổi, dùng $Q=mL$. Thế năng trọng trường của khối lượng $m$ ở độ cao $h$ là $E_p=mgh$; lấy $g=9{,}81\,\text{m}\!\cdot\!\text{s}^{-2}$.
\end{indication}
\begin{solution}
\textbf{1.}
\[
Q_1=mc_{\text{ice}}\Delta T
=1{,}0\times2{,}1\times10^3\times100=2{,}1\times10^5\,\text{J}.
\]
\textbf{2.} Trong khi nóng chảy, nhiệt độ giữ ở $0\,{}^\circ\text{C}$:
\[
Q_2=mL_f=1{,}0\times3{,}34\times10^5=3{,}34\times10^5\,\text{J}.
\]
\textbf{3.}
\[
Q_3=mc_{\text{water}}\Delta T
=1{,}0\times4{,}18\times10^3\times100=4{,}18\times10^5\,\text{J}.
\]
\textbf{4.}
\[
Q_4=mL_v=1{,}0\times2{,}26\times10^6=2{,}26\times10^6\,\text{J}.
\]
Riêng giai đoạn này cần năng lượng cỡ vài megajoule.

\textbf{5.} Hóa hơi là giai đoạn tốn năng lượng nhất, vượt xa các giai đoạn khác:
\[
Q_4=2{,}26\times10^6>Q_3=4{,}18\times10^5>Q_2=3{,}34\times10^5>Q_1=2{,}1\times10^5\,\text{J}.
\]
Hóa hơi cần khoảng $(2{,}26\times10^6)/(4{,}18\times10^5)\approx5{,}4$ lần năng lượng đun nước lỏng từ $0$ lên $100\,{}^\circ\text{C}$.

Với khối lượng $m$ rơi từ $h=10\,\text{m}$, $E_p=mgh$. Đặt $mgh=Q_3$:
\[
m=\frac{Q_3}{gh}=\frac{4{,}18\times10^5}{9{,}81\times10}
\approx4{,}26\times10^3\,\text{kg}.
\]
Vì vậy phải cho khoảng $4{,}3$ tấn rơi mười mét mới giải phóng năng lượng bằng năng lượng đun một kilôgam nước từ $0$ lên $100\,{}^\circ\text{C}$. Giá trị rất lớn này giải thích vì sao đun nước chiếm phần đáng kể trong tiêu thụ năng lượng gia đình. Nhiệt dung riêng của nước cũng rất lớn so với kim loại thông dụng, chẳng hạn gấp khoảng mười lần đồng (xem các bài sau).

\textbf{6.} Ngưng tụ là quá trình ngược của hóa hơi; hơi truyền nhiệt ẩn cho kính:
\[
Q_{\text{rel}}=m_vL_v.
\]
Với $m_v=1{,}0\times10^{-2}\,\text{kg}$,
\[
Q_{\text{rel}}=1{,}0\times10^{-2}\times2{,}45\times10^6
=2{,}45\times10^4\,\text{J}.
\]
Do đó, ngay cả một lượng hơi nhỏ ngưng tụ cũng có thể tỏa nhiệt đáng kể, được kính và môi trường xung quanh hấp thụ.
\end{solution}
\end{exo}

\begin{exo}[Sự thoát hơi nước của cây lớn: máy điều hòa tự nhiên?]{odg-evapotranspiration-arbre}
\seoready{true}
\keywords{bậc độ lớn, thoát hơi nước, cây, nhiệt ẩn, công suất lạnh, máy điều hòa, COP}

Trong một ngày nóng và khô, cây lớn được cung cấp đủ nước có thể thoát hơi khoảng $500\,\text{L}=5{,}0\times10^{-1}\,\text{m}^3$ nước. Vì thoát hơi chủ yếu xảy ra khi có ánh sáng, giả sử quá trình kéo dài mười hai giờ. Tán cây phủ diện tích mặt đất $A_{\text{tree}}=50\,\text{m}^2$. Khối lượng riêng của nước và nhiệt ẩn hóa hơi ở áp suất khí quyển là
\[
\rho_{\text{water}}=1{,}0\times10^3\,\text{kg}\!\cdot\!\text{m}^{-3},
\qquad L_v=2{,}45\times10^6\,\text{J}\!\cdot\!\text{kg}^{-1}.
\]
\begin{enumerate}
  \item Tính khối lượng nước thoát hơi trong ngày và năng lượng cần để hóa hơi lượng nước đó.
  \item Giải thích vì sao quá trình này làm mát.
  \item Tính công suất làm mát trung bình trên mỗi mét vuông trong mười hai giờ thoát hơi tích cực.
  \item Để so sánh, xét máy điều hòa có công suất điện $P_{\mathrm{el}}=2{,}0\times10^3\,\text{W}$, hệ số hiệu suất lạnh $\mathrm{COP}=3{,}0$, làm mát phòng diện tích $A_{\text{room}}=20\,\text{m}^2$. Theo định nghĩa, $\mathrm{COP}=P_{\text{cool}}/P_{\mathrm{el}}$. Tính công suất lạnh và công suất lạnh trên mỗi mét vuông rồi so sánh với cây.
  \item Tại sao không thể kết luận cây và máy điều hòa tạo cảm giác mát giống hệt nhau, dù công suất trên đơn vị diện tích cùng bậc độ lớn?
  \item Có thể làm mát nhà bằng nhiều cây trong nhà không? (Kiến thức chung: câu trả lời phụ thuộc vào các quá trình sinh học.)
\end{enumerate}
\begin{indication}
Dùng $m=\rho V$, $Q_{\mathrm{evap}}=mL_v$, $P=Q/\tau$. Với máy điều hòa, $P_{\text{cool}}=\mathrm{COP}\,P_{\mathrm{el}}$.
\end{indication}
\begin{solution}
\textbf{1.}
\[
m=\rho_{\text{water}}V=1{,}0\times10^3\times5{,}0\times10^{-1}
=5{,}0\times10^2\,\text{kg},
\]
\[
Q_{\mathrm{evap}}=mL_v=5{,}0\times10^2\times2{,}45\times10^6
=1{,}225\times10^9\,\text{J}.
\]
Đó là năng lượng cỡ một gigajoule.

\textbf{2.} Hóa hơi là chuyển pha thu nhiệt. Năng lượng được lấy từ lá, đất và không khí xung quanh, làm giảm nhiệt độ của lá và vùng lân cận.

\textbf{3.} Mười hai giờ là $\tau=12\times3600=4{,}32\times10^4\,\text{s}$. Do đó
\[
P_{\text{tree}}=\frac{Q_{\mathrm{evap}}}{\tau}
=\frac{1{,}225\times10^9}{4{,}32\times10^4}\approx2{,}84\times10^4\,\text{W},
\]
và trên đơn vị diện tích
\[
\frac{P_{\text{tree}}}{A_{\text{tree}}}
=\frac{2{,}84\times10^4}{50}\approx5{,}7\times10^2\,\text{W}\!\cdot\!\text{m}^{-2}.
\]
\textbf{4.}
\[
P_{\text{cool}}=\mathrm{COP}\,P_{\mathrm{el}}
=3{,}0\times2{,}0\times10^3=6{,}0\times10^3\,\text{W},
\]
\[
\frac{P_{\text{cool}}}{A_{\text{room}}}
=\frac{6{,}0\times10^3}{20}=3{,}0\times10^2\,\text{W}\!\cdot\!\text{m}^{-2}.
\]
Trong mô hình này, công suất thoát hơi trên đơn vị diện tích của cây gần gấp đôi máy điều hòa thông thường.

\textbf{5.} So sánh chỉ liên quan đến dòng năng lượng. Máy điều hòa làm mát thể tích kín và được kiểm soát; hơi từ cây phân tán vào khí quyển và gió mang không khí nóng tới. Cảm giác phụ thuộc vào thông gió, độ ẩm, nhiệt độ và khoảng cách đến cây. Sự thoát hơi còn thay đổi theo ánh nắng, gió, độ ẩm, loài cây và lượng nước, và có thể giảm mạnh khi hạn hán. Cây cũng tạo bóng râm, trực tiếp giảm bức xạ mặt trời tới đất và con người. Vì vậy phép tính so sánh bậc độ lớn chứ không chứng minh sự tương đương chính xác về tiện nghi nhiệt.

\textbf{6.} Trên thực tế, tác dụng làm mát của cây trong nhà rất nhỏ. Ở hầu hết cây, ánh sáng thúc đẩy khí khổng mở; hơi nước thoát qua các lỗ nhỏ này trên lá. Ánh sáng trong nhà thường yếu nên khí khổng mở ít và thoát hơi giảm. Trong phòng ít thông gió, độ ẩm tăng cũng dần làm chậm bay hơi.

Nhiều cây có thể gây làm mát bay hơi nhẹ tại chỗ nhưng không thay thế máy điều hòa. Muốn liên tục đưa năng lượng ra khỏi nhà phải thải không khí ẩm ra ngoài. Ánh sáng nhân tạo mạnh có thể kích thích thoát hơi, nhưng điện năng của nó gần như hoàn toàn trở thành nhiệt trong phòng. Thực vật CAM thích nghi với môi trường khô, gồm xương rồng, là ngoại lệ: khí khổng chủ yếu mở ban đêm.
\end{solution}
\end{exo}

\begin{exo}[Trộn hai khối lượng nước]{calorimetrie-melange-eau}
\seoready{true}
\keywords{phép đo nhiệt lượng, trộn, nhiệt dung, nhiệt dung riêng, Joseph Black, nhiệt độ cân bằng}

Đổ $m_1=0{,}200\,\text{kg}$ nước ở $T_1=80\,^\circ\text{C}$ vào bình đã chứa $m_2=0{,}300\,\text{kg}$ nước ở $T_2=15\,^\circ\text{C}$. Bình là nhiệt lượng kế lý tưởng: bỏ qua thất thoát nhiệt ra ngoài và nhiệt dung của bình. Ta có $Q=mc\Delta T$, với $c=4{,}18\times10^3\,\text{J}\!\cdot\!\text{kg}^{-1}\!\cdot\!\text{K}^{-1}$ là nhiệt dung riêng của nước lỏng.
\begin{enumerate}
  \item Giải thích vì sao nhiệt do nước nóng tỏa ra được nước lạnh nhận hoàn toàn và viết phương trình bảo toàn chứa $m_1$, $m_2$, $c$, $T_1$, $T_2$, $T_{eq}$.
  \item Suy ra biểu thức và giá trị số của $T_{eq}$.
  \item Tính nhiệt $Q$ do nước nóng tỏa ra khi trộn.
  \item Có thể dùng hỗn hợp cùng một chất để xác định $c$ không? Giải thích.
\end{enumerate}
\begin{indication}
Nhiệt lượng kế cách nhiệt và cứng nên tổng nội năng $U_1+U_2$ bảo toàn: nhiệt một phần tỏa ra được phần kia nhận với dấu ngược lại.
\end{indication}
\begin{solution}
\textbf{1.}
\[
m_1c(T_{eq}-T_1)+m_2c(T_{eq}-T_2)=0.
\]
\textbf{2.} Hệ số $c$ triệt tiêu:
\[
T_{eq}=\frac{m_1T_1+m_2T_2}{m_1+m_2}
=\frac{0{,}200\times80+0{,}300\times15}{0{,}200+0{,}300}=41\,^\circ\text{C}.
\]
\textbf{3.}
\[
Q=m_1c(T_1-T_{eq})=0{,}200\times4{,}18\times10^3\times(80-41)
\approx3{,}26\times10^4\,\text{J}.
\]
Nước lạnh nhận đúng lượng nhiệt đó:
\[
m_2c(T_{eq}-T_2)=0{,}300\times4{,}18\times10^3\times26
\approx3{,}26\times10^4\,\text{J}.
\]
\textbf{4.} Không. Với cùng một chất, $c$ nhân cả hai số hạng của cân bằng và triệt tiêu. Nhiệt độ cân bằng không phụ thuộc $c$, mà chỉ là trung bình nhiệt độ ban đầu có trọng số khối lượng. Đo nhiệt dung riêng bằng phương pháp trộn cần một chất có nhiệt dung đã biết, như bài sau.
\end{solution}
\end{exo}

\begin{exo}[Xác định nhiệt dung riêng của kim loại bằng phương pháp trộn]{calorimetrie-methode-melanges}
\seoready{true}
\keywords{phép đo nhiệt lượng, phương pháp trộn, nhiệt dung riêng, nhiệt lượng kế, đương lượng nước}

Như Joseph Black đã làm ở thế kỷ XVIII, ta đo nhiệt dung riêng $c_m$ của kim loại chưa biết bằng \emph{phương pháp trộn}. Mẫu có $m=0{,}150\,\text{kg}$ được nung tới $T_m=95\,^\circ\text{C}$ rồi nhanh chóng nhúng vào nhiệt lượng kế chứa $m_e=0{,}250\,\text{kg}$ nước, với $c_e=4{,}18\times10^3\,\text{J}\!\cdot\!\text{kg}^{-1}\!\cdot\!\text{K}^{-1}$ và $T_e=18\,^\circ\text{C}$. Nhiệt độ cân bằng đo được là $T_{eq}=22\,^\circ\text{C}$. Ban đầu bỏ qua nhiệt dung của bình, que khuấy và nhiệt kế.
\begin{enumerate}
  \item Viết bảo toàn năng lượng cho hệ cô lập kim loại–nước, với $c_m$ là ẩn duy nhất.
  \item Suy ra và tính $c_m$. Giá trị gần kim loại thông dụng nào? Các giá trị tham khảo, đơn vị $\text{J}\!\cdot\!\text{kg}^{-1}\!\cdot\!\text{K}^{-1}$: nhôm $9{,}0\times10^2$, sắt $4{,}5\times10^2$, đồng $3{,}9\times10^2$, chì $1{,}3\times10^2$.
  \item Bình cũng hấp thụ nhiệt. Mô hình hóa bằng \emph{đương lượng nước} $\mu=1{,}5\times10^{-2}\,\text{kg}$: nhiệt lượng kế như khối lượng nước bổ sung $\mu$. Tính lại $c_m$ và giải thích chiều hiệu chỉnh.
  \item Vì sao phải nhúng mẫu nhanh và cách nhiệt tốt nhiệt lượng kế với không khí xung quanh?
\end{enumerate}
\begin{indication}
Kim loại tỏa $Q_m=mc_m(T_m-T_{eq})$, được nước và, ở câu 3, nhiệt lượng kế nhận hoàn toàn: $Q_m=(m_e+\mu)c_e(T_{eq}-T_e)$.
\end{indication}
\begin{solution}
\textbf{1.}
\[
mc_m(T_m-T_{eq})=m_ec_e(T_{eq}-T_e).
\]
\textbf{2.}
\[
c_m=\frac{m_ec_e(T_{eq}-T_e)}{m(T_m-T_{eq})}
=\frac{0{,}250\times4{,}18\times10^3\times(22-18)}{0{,}150\times(95-22)}
\approx3{,}82\times10^2\,\text{J}\!\cdot\!\text{kg}^{-1}\!\cdot\!\text{K}^{-1}.
\]
Giá trị này rất gần với đồng.

\textbf{3.}
\[
c_m=\frac{(m_e+\mu)c_e(T_{eq}-T_e)}{m(T_m-T_{eq})}
=\frac{(0{,}250+0{,}015)\times4{,}18\times10^3\times4}{0{,}150\times73}
\approx4{,}05\times10^2\,\text{J}\!\cdot\!\text{kg}^{-1}\!\cdot\!\text{K}^{-1}.
\]
Hiệu chỉnh làm $c_m$ tăng: bỏ qua bình đã đánh giá thấp nhiệt thực sự do kim loại tỏa ra và do đó đánh giá thấp nhiệt dung riêng.

\textbf{4.} Phương pháp giả thiết hệ cô lập suốt phép đo. Nhúng nhanh hạn chế trao đổi nhiệt giữa mẫu nóng và không khí trước khi mẫu vào nước; cách nhiệt tốt hạn chế mất nhiệt khi cân bằng hình thành. Mọi truyền nhiệt không tính đến đều làm sai cân bằng và $c_m$. Đây chính là sự cẩn trọng thực nghiệm của Black, rồi Lavoisier và Laplace với nhiệt lượng kế băng.
\end{solution}
\end{exo}

\begin{exo}[Calo thực phẩm và công suất trao đổi chất]{calorimetrie-calories-alimentaires}
\seoready{true}
\keywords{calo, kilocalo, Calo thực phẩm, đương lượng cơ học, công suất trao đổi chất, đơn vị}

Nhãn dinh dưỡng cho biết một người lớn tiêu thụ trung bình khoảng $2000\,\text{kcal}$ mỗi ngày. Kilocalo là đơn vị ngoài SI vẫn dùng trong dinh dưỡng; $1\,\text{kcal}=4{,}18\times10^3\,\text{J}$.
\begin{enumerate}
  \item Đổi năng lượng hằng ngày sang joule.
  \item Suy ra công suất trao đổi chất trung bình bằng watt, giả sử năng lượng được dùng đều trong 24 giờ.
  \item Một thanh năng lượng ghi ``$210\,\text{kcal}$''. Nếu toàn bộ năng lượng chuyển thành nhiệt, có thể đun bao nhiêu nước từ $20\,^\circ\text{C}$ đến sôi ở $100\,^\circ\text{C}$? Lấy $c_{\text{water}}=4{,}18\times10^3\,\text{J}\!\cdot\!\text{kg}^{-1}\!\cdot\!\text{K}^{-1}$.
  \item Cơ thể người thực sự sử dụng năng lượng này dưới những dạng nào? Trả lời định tính.
\end{enumerate}
\begin{indication}
Câu 2 dùng $\mathcal P=E/\Delta t$. Câu 3 đổi sang joule rồi giải $m$ từ $Q=mc\Delta T$.
\end{indication}
\begin{solution}
\textbf{1.}
\[
E=2000\times4{,}18\times10^3=8{,}36\times10^6\,\text{J}.
\]
\textbf{2.}
\[
\mathcal P=\frac{8{,}36\times10^6}{8{,}64\times10^4}\approx97\,\text{W}.
\]
Trong cả ngày, một người dùng công suất trung bình gần bằng bóng đèn sợi đốt. Bậc độ lớn thường gây ngạc nhiên này cho thấy ích lợi của việc đổi calo thực phẩm sang đơn vị vật lý quen thuộc.

\textbf{3.}
\[
Q=210\times4{,}18\times10^3\approx8{,}78\times10^5\,\text{J}.
\]
Với $\Delta T=80\,\text{K}$,
\[
m=\frac{Q}{c\Delta T}
=\frac{8{,}78\times10^5}{4{,}18\times10^3\times80}
\approx2{,}63\,\text{kg}.
\]
Theo bậc độ lớn, một thanh năng lượng đủ để đun sôi khoảng $2{,}5\,\text{kg}$ nước máy.

\textbf{4.} Cơ thể không “tiêu thụ” năng lượng theo nghĩa làm nó biến mất, mà chuyển hóa năng lượng hóa học của chất dinh dưỡng. Một phần dùng cho công cơ, hoạt động cơ quan, duy trì gradient điện của tế bào, vận chuyển phân tử và tổng hợp hóa học; một phần có thể tích trữ dưới dạng glycogen hoặc mỡ. Cuối cùng, phần lớn năng lượng đã dùng bị tiêu tán dưới dạng nhiệt, giúp duy trì thân nhiệt trước khi truyền ra môi trường. Một phần nhỏ rời cơ thể dưới dạng năng lượng hóa học còn trong chất thải.
\end{solution}
\end{exo}
